\hypertarget{chapter-08-page-169}{}
\subsection{若干向量值不等式}
\label{sec:ch8-vector-valued-inequalities}

我们将通过第~\ref{sec:ch5-vector-valued-extension}~节引入的向量值奇异积分理论研究 Littlewood--Paley 理论. 本节利用这些思想以及第~\ref{chap:weighted-inequalities}~章的加权范数不等式, 证明后文所需的若干向量值不等式.

\begin{Theorem}
\label{thm:ch8-vector-valued-convolution}
设 $T$ 是在 $L^2(\mathbb R^n)$ 上有界的卷积算子, 其相伴核 $K$ 满足 Hörmander 条件 \eqref{eq:ch5-hormander-condition}. 则对任意 $1<r<\infty$ 和 $1<p<\infty$,
\[
 \left\|\left(\sum_j|Tf_j|^r\right)^{1/r}\right\|_p
 \le C_{p,r}\left\|\left(\sum_j|f_j|^r\right)^{1/r}\right\|_p;
\]
当 $p=1$ 时,
\[
 \left|\left\{x\in\mathbb R^n:
 \left(\sum_j|Tf_j(x)|^r\right)^{1/r}>\lambda\right\}\right|
 \le\frac{C_r}{\lambda}
 \left\|\left(\sum_j|f_j|^r\right)^{1/r}\right\|_1.
\]
\end{Theorem}

\begin{Proof}
在定理~\ref{thm:ch5-vector-valued-calderon-zygmund} 中取 $A=B=\ell^r$. 当 $p=r$ 时, 第一个不等式直接来自 $T$ 在 $L^r(\mathbb R^n)$ 上的有界性, $1<r<\infty$. 现在考虑把序列 $\{f_j\}$ 映为 $\{Tf_j\}$ 的向量值算子. 其核为 $K(x)I$, 其中 $I$ 是 $\ell^r$ 上的恒等算子. 对卷积算子, 条件 \eqref{eq:ch5-vector-hormander-y} 和 \eqref{eq:ch5-vector-hormander-x} 都化为
\[
 \int_{|x|>2|y|}\|K(x-y)I-K(x)I\|_{\ell^r}\,dx\le C,
\]
而
\[
 \|K(x-y)I-K(x)I\|_{\ell^r}=|K(x-y)-K(x)|,
\]
所以该条件正是 Hörmander 条件.
\end{Proof}

\hypertarget{chapter-08-page-170}{}
一个自然问题是把定理~\ref{thm:ch8-vector-valued-convolution} 推广到算子序列 $\{T_j\}$, 得到
\begin{equation}
\label{eq:ch8-vector-operator-sequence}
 \left\|\left(\sum_j|T_jf_j|^r\right)^{1/r}\right\|_p
 \le C_{p,r}\left\|\left(\sum_j|f_j|^r\right)^{1/r}\right\|_p,
 \qquad1<p,r<\infty.
\end{equation}
似乎可以猜想: 若 $T_j$ 在 $L^2$ 上一致有界, 且其核 $K_j$ 一致满足 Hörmander 条件, 则 \eqref{eq:ch8-vector-operator-sequence} 成立. 然而当 $p\ne r$ 时, 这仍可能不成立. 把 $\{f_j\}$ 映为 $\{T_jf_j\}$ 的向量值算子, 其核是 $\mathcal L(\ell^r,\ell^r)$ 中把 $\{\lambda_j\}$ 映为 $\{K_j(x)\lambda_j\}$ 的算子. 对该算子, Hörmander 条件变为
\[
 \int_{|x|>2|y|}\sup_j|K_j(x-y)-K_j(x)|\,dx\le C,
\]
而这一条件足以保证 \eqref{eq:ch8-vector-operator-sequence}.

不过, 定理~\ref{thm:ch8-vector-valued-convolution} 可在特殊情形下给出 \eqref{eq:ch8-vector-operator-sequence} 一类不等式.

\begin{Corollary}
\label{cor:ch8-interval-projections-vector-valued}
设 $\{I_j\}$ 是实直线上由有限或无限区间组成的序列, 算子 $S_j$ 由
\[
 \widehat{S_jf}(\xi)=\chi_{I_j}(\xi)\widehat f(\xi)
\]
定义. 则对 $1<r,p<\infty$,
\[
 \left\|\left(\sum_j|S_jf_j|^r\right)^{1/r}\right\|_p
 \le C_{p,r}\left\|\left(\sum_j|f_j|^r\right)^{1/r}\right\|_p.
\]
\end{Corollary}

\begin{Proof}
若 $I_j=(a_j,b_j)$, 则由 Hilbert 变换的 Fourier 乘子表示,
\[
 S_jf_j=\frac i2\left(M_{a_j}HM_{-a_j}f_j-M_{b_j}HM_{-b_j}f_j\right),
\]
区间无界时作显然修改. 对 Hilbert 变换应用定理~\ref{thm:ch8-vector-valued-convolution}, 即得结论.
\end{Proof}

证明 \eqref{eq:ch8-vector-operator-sequence} 一类向量值不等式的常用方法涉及加权范数不等式. 下一结果是这种方法的一个例子.

\hypertarget{chapter-08-page-171}{}
\begin{Theorem}
\label{thm:ch8-weighted-sequence-square}
设 $\{T_j\}$ 是线性算子序列. 对任意 $w\in A_2$, 每个 $T_j$ 都在 $L^2(w)$ 上有界, 常数关于 $j$ 一致且只依赖于 $w$ 的 $A_2$ 常数. 则对所有 $1<p<\infty$,
\begin{equation}
\label{eq:ch8-weighted-sequence-square}
 \left\|\left(\sum_j|T_jf_j|^2\right)^{1/2}\right\|_p
 \le C\left\|\left(\sum_j|f_j|^2\right)^{1/2}\right\|_p.
\end{equation}
\end{Theorem}

\begin{Proof}
当 $p=2$ 时结论直接成立. 设 $p>2$. 存在范数为 $1$ 的 $u\in L^{(p/2)'}$, 使
\[
 \left\|\left(\sum_j|T_jf_j|^2\right)^{1/2}\right\|_p^2
 =\int_{\mathbb R^n}\sum_j|T_jf_j|^2u.
\]
由定理~\ref{thm:ch7-a1-characterization}, 若 $0<\delta<1$, 则 $M(u^{1/\delta})^\delta$ 是 $A_1$ 权, 因而也属于 $A_2$, 且其常数只依赖于 $\delta$. 又因几乎处处有 $u(x)\le M(u^{1/\delta})(x)^\delta$,
\[
 \left\|\left(\sum_j|T_jf_j|^2\right)^{1/2}\right\|_p^2
 \le C_\delta\int_{\mathbb R^n}\sum_j|f_j|^2M(u^{1/\delta})^\delta.
\]
取 $\delta$ 使 $\delta(p/2)'>1$, 再用指数 $p/2$ 和 $(p/2)'$ 的 Hölder 不等式, 得 \eqref{eq:ch8-weighted-sequence-square}. 最后, 当 $p<2$ 时, 伴随算子 $T_j^*$ 对 $w\in A_2$ 也在 $L^2(w)$ 上有界, 所以由对偶性得到结论.
\end{Proof}
